\section{A kernel theorem for recurrent nonuniform transport}\label{sec:kernel}
The moment condition in the transfer theorem is useful only insofar as it can be derived without solving the approximation problem itself. We now derive it from a positive operator attached to a marked factor of the raw experiment. The factor may be used only for analysis. It need not be observed, and it is never supplied to the encoder.

\subsection{One-step data and positive resolvents}
Let $\Xi$ be standard Borel, and suppose that the coefficients of the innovation recurrence are generated by a time-homogeneous marked kernel
\begin{equation}\label{eq:marked-kernel}
 P(d\xi',d\gamma,d\mathbf b\mid\xi),\qquad
 \gamma\ge0,\quad \mathbf b\in[0,\infty)^J,\quad \sum_j b_j\le B_*.
\end{equation}
Thus, conditionally on $\xi_t$, the next environment and marks are independent of all previous marks and of the seed. This assertion is checked from the raw preparation, the chosen exploration and the one-step update bounds. It excludes coefficients chosen after seeing the grid allocation. Controller state may have to be included to make the factor Markov; its implementation cost is not thereby removed.

Assume that the environment marginal has an invariant probability $\mu$. Disintegrate the stationary marked edge measure in the reverse direction:
\[
 \mu(d\xi)P(d\xi',d\gamma,d\mathbf b\mid\xi)
   =\mu(d\xi')\overleftarrow P(d\xi,d\gamma,d\mathbf b\mid\xi').
\]
All the following identities are statements about measurable versions $\mu$-almost everywhere. Define, for nonnegative measurable $f$,
\begin{align}
 (K_p f)(\xi')&=\int\gamma^p f(\xi)\,d\overleftarrow P,
       &&p=1,2,\label{eq:positive-operators}\\
 f_j(\xi')&=\int b_j\,d\overleftarrow P,
 &g_{jk}(\xi')&=\int b_jb_k\,d\overleftarrow P,\notag\\
 (D_jf)(\xi')&=\int\gamma b_j f(\xi)\,d\overleftarrow P.\notag
\end{align}
The structural hypothesis is the one-step inequality
\begin{equation}\label{eq:kernel-drift}
 1\le v\le v_+<\infty,\qquad K_2v\le\lambda v,
 \qquad 0\le\lambda<1.
\end{equation}
Here $v$ is a measurable Lyapunov weight. Neither $\gamma<1$ on each edge nor a bound on a previously unrolled error is assumed. The operators depend only on the experimental factor and the selected local update/rounding bounds.

For a positive weight $u$, write $\|f\|_u=\operatorname*{ess\,sup}|f|/u$. Cauchy--Schwarz gives $K_1\sqrt v\le\sqrt\lambda\sqrt v$. Hence the positive Neumann resolvents $\mathcal S_1=(I-K_1)^{-1}$ on the $\sqrt v$-weighted space and $\mathcal S_2=(I-K_2)^{-1}$ on the $v$-weighted space exist. Set
\begin{align}
 h_j&=\mathcal S_1 f_j,\label{eq:resolvent-profiles}\\
 V_{jk}&=\mathcal S_2\bigl(g_{jk}+D_jh_k+D_kh_j\bigr),
 &H_{jk}&=\int V_{jk}\,d\mu.\notag
\end{align}
The profiles are finite. Indeed $h_j$ is bounded by a constant times $\sqrt v$, and Cauchy--Schwarz bounds $D_jh_k$ by a constant times $\sqrt v$, hence by a constant times $v$. The constants depend only on $B_*,v_+,1-\lambda$ and $J$, not on a memory allocation.

\begin{theorem}[Forced moment resolvent from a raw-kernel factor]\label{thm:kernel}
Suppose \eqref{eq:marked-kernel}--\eqref{eq:kernel-drift} hold. Let
\[
 Z_{0,j}=0,\qquad Z_{t+1,j}=\gamma_{t+1}Z_{t,j}+b_{t+1,j}.
\]
Under the stationary environment preparation, for every $N\ge0$ and every $z\in[0,\infty)^J$,
\begin{equation}\label{eq:kernel-energy}
 \E\left(\sum_j z_jZ_{N,j}\right)^2\le z^{\mathsf T}Hz.
\end{equation}
The matrix $H$ is symmetric positive semidefinite and has nonnegative entries. The same inequality holds with a nonnegative seed if, conditionally on $\xi_0$, its first moments are bounded by $h_j$ and its mixed second moments by $V_{jk}$, entry by entry. If the actual initial environment law has density at most $D_*$ with respect to $\mu$, and uses the same conditional seed rule, the right side is multiplied by $D_*$.

Consequently, for the actual acquired weights $w_{N,j}$, condition \textup{(T)} follows with
\begin{equation}\label{eq:resolvent-gamma}
 \Gamma_N^{\rm ker}
   =D_*\left\|\operatorname{diag}(w_N)^{-1/2}
                    H\operatorname{diag}(w_N)^{-1/2}\right\|_{\rm op},
\end{equation}
provided every zero-weight row of $H$ is zero; the inverse is taken only on positive-weight coordinates. If this proviso fails, this stationary majorant does not certify that horizon. A finite-horizon profile or another mechanism must then be used. In particular, for $J=1$, $w_{N,1}=1$, the one-step drift alone gives a uniform finite transport constant, with no assumption of uniform contraction of the active updates.
\end{theorem}
\begin{proof}
First start the environment at $\mu$ and the forcing process at zero. Define conditional moment profiles
\[
 m_{t,j}(\xi)=\E[Z_{t,j}\mid\xi_t=\xi],\qquad
 C_{t,jk}(\xi)=\E[Z_{t,j}Z_{t,k}\mid\xi_t=\xi].
\]
The Markov-factor hypothesis ensures conditional independence of future edge marks and the previous forcing history given the previous environment. Conditioning on the stationary incoming edge therefore gives the exact recursions
\begin{align}
 m_{t+1,j}&=K_1m_{t,j}+f_j,\label{eq:profile-recursion}\\
 C_{t+1,jk}&=K_2C_{t,jk}+D_jm_{t,k}+D_km_{t,j}+g_{jk}.\notag
\end{align}
Truncation of the marks justifies these identities before integrability is known; monotone convergence then removes the truncation. Alternatively, induction using the finite bounds below proves integrability at each finite time.

Positivity and the identities $h_j=K_1h_j+f_j$ and
$V_{jk}=K_2V_{jk}+D_jh_k+D_kh_j+g_{jk}$ imply
$0\le m_{t,j}\le h_j$ and $0\le C_{t,jk}\le V_{jk}$ by induction. In fact $m_t$ and $C_t$ increase from zero. The first recursion converges in the $\sqrt v$ norm to $h$. Subtracting the second recursion from its fixed-point identity, or summing its geometric convolution, proves convergence in the $v$ norm to $V$. To see the latter directly, split the convolution at a fixed lag: its long-lag part is bounded by a geometric $\lambda$-tail, and the finitely many remaining terms converge because $m_t\to h$.

For each $t$ the conditional matrix $C_t(\xi)$ is positive semidefinite. Entrywise convergence in a finite matrix preserves this property outside one common null set. Thus $V(\xi)$ and its integral $H$ are positive semidefinite. All entries are nonnegative. Integrating the entrywise upper bounds and multiplying by $z_jz_k\ge0$ proves \eqref{eq:kernel-energy}. This last nonnegative-cone bound is the one needed for covering radii; it is not an entrywise-to-Loewner inference.

For a seed satisfying the stated conditional inequalities, the same induction in \eqref{eq:profile-recursion} applies, without the monotone-convergence assertion for that seed. An initial density bounded by $D_*$ dominates the entire path-and-seed law by $D_*$, since all subsequent kernels agree. Integrating any nonnegative path functional proves the domination claim. Finally apply the operator norm bound to $\operatorname{diag}(w_N)^{1/2}z$. A zero diagonal of a positive semidefinite matrix forces its row to be zero; the explicit proviso prevents assigning positive transport energy to a component of zero acquired probability.
\end{proof}

This theorem replaces a path-length family of forced-moment inequalities by one-step kernel integration, a Lyapunov inequality, and two positive resolvents. The comparison with actual acquired weights remains visible in \eqref{eq:resolvent-gamma}: small mass is not made free by a stability theorem. Uniform matching across a family means that this explicitly computed normalized matrix is uniformly bounded. For one acquired chart no further mass-normalization condition is needed. For a finite environmental chain, all these objects are finite matrices and vectors, independent of $N$ and $m$.

\begin{corollary}[Finite-state check and recurrent expansion]\label{cor:finite-kernel}
For a finite environment, $\rho(K_2)<1$ is equivalent to the existence of a strictly positive vector $v$ and $\lambda<1$ satisfying $K_2v\le\lambda v$. In particular consider a two-state chain, with states $C,E$, transition matrix
\[
 P=\begin{pmatrix}1-p&p\\1&0\end{pmatrix},\qquad 0<p\le1,
\]
and gain $1/4$ on arrival at $C$, $2$ on arrival at $E$. Under its invariant law, the reversed chain has the same transition matrix and
\[
 K_2=\begin{pmatrix}(1-p)/16&p/16\\4&0\end{pmatrix},
 \qquad K_2\binom{1}{8}\le\tfrac12\binom{1}{8}.
\]
Bounded forcing therefore satisfies Theorem~\ref{thm:kernel}, although expansions recur and the conditional squared gain on arrival at $E$ is four.
\end{corollary}
\begin{proof}
A positive vector inequality makes the weighted sup norm of $K_2$ at most $\lambda$, so its spectral radius is less than one. Conversely, when $\rho(K_2)<1$, let $v=(I-K_2)^{-1}\mathbf1\ge\mathbf1$. Then $K_2v=v-\mathbf1\le(1-1/\max_i v_i)v$; the zero matrix case is included. The displayed chain has invariant probabilities $(1,p)/(1+p)$ and satisfies detailed balance. Multiplication by the squared incoming gain gives the stated matrix. Its first row applied to $(1,8)$ is $(1+7p)/16\le1/2$ and its second is four. These calculations prove the example. It is an operator witness; a complete raw experiment and its acquired geometry are verified separately below.
\end{proof}

\begin{proposition}[Exact suspension preserves the forced profiles]\label{prop:suspension}
Replace each transition of \eqref{eq:marked-kernel} independently by an exact hold with probability $1-\theta$, $0<\theta\le1$. A hold has unchanged environment, gain one and zero forcing. Keep every hold in the raw-call and physical-time accounts. The invariant environment law remains $\mu$, and the profiles $h,V,H$ in \eqref{eq:resolvent-profiles} are unchanged. Thus the bound \eqref{eq:resolvent-gamma} has no inverse power of $\theta$ apart from any change in the independently computed acquired weights. At $\theta=0$ the forcing process remains its seed, so the same upper bound holds for the seeds admitted in Theorem~\ref{thm:kernel}.
\end{proposition}
\begin{proof}
Reversing the stationary held kernel gives the same mixture. Therefore
\[
 K_p^\theta=(1-\theta)I+\theta K_p,
 \quad f_j^\theta=\theta f_j,\quad
 g_{jk}^\theta=\theta g_{jk},\quad D_j^\theta=\theta D_j.
\]
The drift constant becomes $1-\theta(1-\lambda)<1$. Cancelling $\theta$ in the two fixed-point equations proves $h^\theta=h$ and $V^\theta=V$. This cancellation is justified by the positive resolvent's uniqueness, not by replacing physical time with a random stopping time. At zero activity there are no new innovations and direct retention proves the assertion. The independence and state-independent value of $\theta$ are part of this proposition; arbitrary adaptive suspensions require their own kernel calculation.
\end{proof}

For a scalar independent active gain $G$, with forcing one at every active step, write $a_1=\E G$ and $a_2=\E G^2<1$. The equations give
\begin{equation}\label{eq:scalar-resolvent}
 h=\frac1{1-a_1},\qquad H=\frac{1+a_1}{(1-a_1)(1-a_2)}.
\end{equation}
Indeed $a_1\le\sqrt{a_2}<1$, and substitution into $H=a_2H+2a_1h+1$ proves the formula. It also bounds any seed with first moment at most $h$ and second moment at most $H$. This explicit second-moment formula is a consequence of the general kernel theorem, not a replacement for it. Section~\ref{sec:recurrent} verifies a recurrent continuous-state experiment in this class, including its actual density and finite program.

A negative average logarithmic gain would not replace \eqref{eq:kernel-drift}. For example, $G=1/4$ with probability $9/10$ and $G=4$ with probability $1/10$ has negative mean logarithm but $\E G^2=53/32>1$. Already the squared product multiplying one earlier nonzero innovation has expectation growing geometrically. The theorem controls forced second moments, not only almost-sure contraction of unforced iterates.
